\begin{pastexam}[label={px:ncku-114-ee-ds-11}]{成大電機所 114 學年度 單選題第 11 題}
（單選；每題 5 分，答錯倒扣 2 分）

Given a hash table with 11 buckets; labeled from 0 to 10. The hash function h(key)= key \% 11
allocates data into a hash table. The chaining strategy is used to store data in each bucket.
Please show the total number of comparisons required to store the numbers ``3, 41, 15, 36, 74,
58, 91, 45, 48, 64''.

a) 20 \quad b) 16 \quad c) 22 \quad d) 14 \quad e) None of the above.
\end{pastexam}

\begin{pastsoln}
答案是 e) None of the above。

題目沒有定義「一次比較」怎麼算。 鏈結法（資結書第~12.2~節）插入時，標準的做法是先在桶子的串列
裡找有沒有相同的鍵值（避免重複），找不到才放進去，所以比較次數＝插入當時串列裡已有的元素個數。
\begin{center}
\begin{tabular}{cccl}
\toprule
鍵值 & $h = \mathit{key} \bmod 11$ & 插入前桶子裡有 & 比較次數 \\
\midrule
3 & 3 & （空） & 0 \\
41 & 8 & （空） & 0 \\
15 & 4 & （空） & 0 \\
36 & 3 & 3 & 1 \\
74 & 8 & 41 & 1 \\
58 & 3 & 3, 36 & 2 \\
91 & 3 & 3, 36, 58 & 3 \\
45 & 1 & （空） & 0 \\
48 & 4 & 15 & 1 \\
64 & 9 & （空） & 0 \\
\bottomrule
\end{tabular}
\end{center}
（$41 = 33 + 8$，$74 = 66 + 8$，$58 = 55 + 3$，$91 = 88 + 3$，$45 = 44 + 1$，$48 = 44 + 4$，
$64 = 55 + 9$。）總共 $1 + 1 + 2 + 3 + 1 = 8$ 次。

另一種常見的算法是每次插入都把「走到串列尾端」也算進去，也就是每筆多 $1$ 次，共 $8 + 10 = 18$ 次。
直接插在串列前端、不檢查重複的話是 $0$ 次。 不論哪一種算法，都得不到 $20$、$16$、$22$、$14$，
所以選 e)。 最後的桶子：$1$：45；$3$：3, 36, 58, 91；$4$：15, 48；$8$：41, 74；$9$：64。
\end{pastsoln}
